\section{Lines in the Cartesian Plane}

A line is one of the simplest geometric subsets of the plane, but it already
shows how several mathematical descriptions can refer to the same object.  We
may draw a line, describe how its coordinates change, write a condition
satisfied by its points, or define it explicitly as a subset of
$\mathbb{R}^2$.  These descriptions should not be confused with one another.

\subsection{A Line as a Set of Points}

Consider the condition
\[
y=2x+1.
\]
By itself, this is a relation between the coordinates $x$ and $y$.  The
geometric object determined by this relation is the set
\[
L
=
\{(x,y)\in\mathbb{R}^2:y=2x+1\}.
\]
We read this as follows:

\begin{quote}
$L$ is the set of all points $(x,y)$ in the Cartesian plane such that
$y=2x+1$.
\end{quote}

Every value of $x$ determines exactly one value of $y$.  For example,
\[
\begin{array}{c|ccccc}
x&-2&-1&0&1&2\\
\hline
y=2x+1&-3&-1&1&3&5
\end{array}
\]
so the points
\[
(-2,-3),\quad(-1,-1),\quad(0,1),\quad(1,3),\quad(2,5)
\]
belong to $L$.

\begin{center}
\begin{tikzpicture}[scale=0.75]
    \draw[axis] (-3.4,0) -- (3.6,0) node[right] {$x$};
    \draw[axis] (0,-4.2) -- (0,6.0) node[above] {$y$};
    \draw[subset,domain=-2.4:2.4,samples=2] plot (\x,{2*\x+1});
    \foreach \x/\y in {-2/-3,-1/-1,0/1,1/3,2/5} {
        \fill[geometricSubsetColor] (\x,\y) circle (2pt);
    }
    \node[subset label,right] at (2.1,5.2) {$L$};
\end{tikzpicture}
\end{center}

The listed points are examples of elements of $L$.  The line itself contains
infinitely many points and is defined by the complete set-builder expression.

\subsection{Slope and Vertical Change}

In the equation
\[
y=mx+b,
\]
the number $m$ measures how the vertical coordinate changes when the horizontal
coordinate changes.  Suppose we increase $x$ by an amount $\Delta x$.  Then
\[
\Delta y=m\,\Delta x.
\]
Therefore
\[
m=\frac{\Delta y}{\Delta x},
\qquad \Delta x\neq 0.
\]
The number $m$ is called the \emph{slope} of the line.

For the line $y=2x+1$, increasing $x$ by $1$ increases $y$ by $2$.  Increasing
$x$ by $3$ increases $y$ by $6$.  The same ratio appears everywhere on the
line:
\[
\frac{\Delta y}{\Delta x}=2.
\]

\begin{center}
\begin{tikzpicture}[scale=0.85]
    \draw[axis] (-0.8,0) -- (4.4,0) node[right] {$x$};
    \draw[axis] (0,-0.8) -- (0,6.0) node[above] {$y$};
    \draw[subset,domain=-0.2:2.5,samples=2] plot (\x,{2*\x+1});
    \fill[geometricSubsetColor] (0.5,2) circle (2pt);
    \fill[geometricSubsetColor] (2,5) circle (2pt);
    \draw[helper] (0.5,2) -- (2,2) node[midway,below] {$\Delta x$};
    \draw[helper] (2,2) -- (2,5) node[midway,right] {$\Delta y$};
\end{tikzpicture}
\end{center}

The sign of the slope describes the direction of the line:
\begin{itemize}
    \item if $m>0$, the line rises as we move from left to right;
    \item if $m<0$, the line falls as we move from left to right;
    \item if $m=0$, the line is horizontal.
\end{itemize}

The slope does not describe one chosen pair of points only.  It is a property of
the entire line.

\subsection{The Meaning of the Constant Term}

In the equation
\[
y=mx+b,
\]
setting $x=0$ gives
\[
y=b.
\]
Thus the line meets the $y$-axis at the point
\[
(0,b).
\]
The number $b$ is called the \emph{$y$-intercept}.

For example, the line
\[
L=\{(x,y)\in\mathbb{R}^2:y=-\tfrac12x+3\}
\]
has slope $-\tfrac12$ and meets the $y$-axis at $(0,3)$.

To find where a non-horizontal line meets the $x$-axis, we set $y=0$.  For the
same line,
\[
0=-\frac12x+3,
\]
so $x=6$.  Hence the line meets the $x$-axis at $(6,0)$.

\subsection{Vertical Lines}

Not every line can be written in the form $y=mx+b$.  A vertical line has the
form
\[
V_c
=
\{(x,y)\in\mathbb{R}^2:x=c\},
\]
where $c$ is a fixed real number.

For example,
\[
V_{-2}
=
\{(x,y)\in\mathbb{R}^2:x=-2\}
\]
is the vertical line through all points whose first coordinate is $-2$.

\begin{center}
\begin{tikzpicture}[scale=0.8]
    \draw[axis] (-4.0,0) -- (3.5,0) node[right] {$x$};
    \draw[axis] (0,-3.2) -- (0,3.5) node[above] {$y$};
    \draw[subset] (-2,-3) -- (-2,3);
    \node[subset label,above] at (-2,3) {$V_{-2}$};
    \draw (-2,0.08) -- (-2,-0.08) node[below=3pt] {$-2$};
\end{tikzpicture}
\end{center}

A vertical line has no finite slope, because moving between two distinct points
on it gives $\Delta x=0$, and the ratio
\[
\frac{\Delta y}{\Delta x}
\]
is then undefined.

\subsection{The General Linear Equation}

A line can also be described by a condition of the form
\[
ax+by=c,
\]
where $a$, $b$, and $c$ are real numbers and $a$ and $b$ are not both zero.  The
corresponding subset of the plane is
\[
L_{a,b,c}
=
\{(x,y)\in\mathbb{R}^2:ax+by=c\}.
\]

When $b\neq 0$, we may solve for $y$:
\[
y=-\frac{a}{b}x+\frac{c}{b}.
\]
Thus the slope is
\[
m=-\frac{a}{b},
\]
and the $y$-intercept is $c/b$.

When $b=0$, the condition becomes
\[
ax=c,
\]
so, because $a\neq 0$,
\[
x=\frac{c}{a}.
\]
This is a vertical line.  The general equation therefore includes both
non-vertical and vertical lines.

\begin{examplebox}
Consider
\[
2x+3y=6.
\]
The geometric object is
\[
L
=
\{(x,y)\in\mathbb{R}^2:2x+3y=6\}.
\]
Solving for $y$ gives
\[
y=-\frac23x+2.
\]
The same line therefore has slope $-\frac23$, meets the $y$-axis at $(0,2)$,
and meets the $x$-axis at $(3,0)$.
\end{examplebox}

\subsection{A Line Through Two Points}

Let
\[
P=(x_1,y_1),
\qquad
Q=(x_2,y_2),
\]
with $x_1\neq x_2$.  The slope of the unique line through these points is
\[
m=\frac{y_2-y_1}{x_2-x_1}.
\]
A point $(x,y)$ lies on the same line precisely when its change from $P$ has the
same slope:
\[
\frac{y-y_1}{x-x_1}=m.
\]
Multiplying by $x-x_1$ gives the point--slope condition
\[
y-y_1=m(x-x_1).
\]
The corresponding line is therefore
\[
L
=
\{(x,y)\in\mathbb{R}^2:y-y_1=m(x-x_1)\}.
\]

If $x_1=x_2$, the two distinct points determine the vertical line
\[
L=\{(x,y)\in\mathbb{R}^2:x=x_1\}.
\]

\begin{examplebox}
The points $P=(1,2)$ and $Q=(4,8)$ determine a line with slope
\[
m=\frac{8-2}{4-1}=2.
\]
Using $P$, its condition is
\[
y-2=2(x-1),
\]
which simplifies to
\[
y=2x.
\]
Hence the line is
\[
L=\{(x,y)\in\mathbb{R}^2:y=2x\}.
\]
\end{examplebox}

\subsection{Parallel Lines: Geometry First, Criterion Second}

Two distinct lines are called \emph{parallel} when they have no common point.
Thus, for subsets $L_1,L_2\subseteq\mathbb{R}^2$, parallelism means
\[
L_1\neq L_2
\qquad\text{and}\qquad
L_1\cap L_2=\varnothing.
\]
This is the geometric definition.  We now ask how it can be recognised from
equations.

Let
\[
L_1=\{(x,y)\in\mathbb{R}^2:y=m_1x+b_1\},
\qquad
L_2=\{(x,y)\in\mathbb{R}^2:y=m_2x+b_2\}.
\]
A common point would have to satisfy both conditions, so
\[
m_1x+b_1=m_2x+b_2.
\]
Equivalently,
\[
(m_1-m_2)x=b_2-b_1.
\]

If $m_1\neq m_2$, this equation has exactly one solution for $x$, and therefore
the two lines meet at exactly one point.  They are not parallel.

If $m_1=m_2$, the equation becomes
\[
0=b_2-b_1.
\]
There are then two possibilities:
\begin{itemize}
    \item if $b_1=b_2$, the two set-builder descriptions define the same line;
    \item if $b_1\neq b_2$, no point satisfies both conditions, so the lines are
    distinct and have empty intersection.
\end{itemize}

We have therefore derived the criterion:
\[
\boxed{L_1\parallel L_2
\iff m_1=m_2\text{ and }b_1\neq b_2.}
\]
The equality of slopes is not the definition of parallelism.  It is an
algebraic consequence of the geometric definition for non-vertical lines.
Two distinct vertical lines are parallel for the same reason: the conditions
$x=c_1$ and $x=c_2$ cannot hold simultaneously when $c_1\neq c_2$.

\begin{center}
\begin{tikzpicture}[scale=0.75]
    \draw[axis] (-4.2,0) -- (4.4,0) node[right] {$x$};
    \draw[axis] (0,-3.5) -- (0,4.0) node[above] {$y$};
    \draw[subset,domain=-3.2:3.2,samples=2] plot (\x,{0.6*\x+1.3});
    \draw[subset,domain=-3.2:3.2,samples=2] plot (\x,{0.6*\x-1.2});
    \node[subset label,right] at (3.0,3.1) {$L_1$};
    \node[subset label,right] at (3.0,0.6) {$L_2$};
\end{tikzpicture}
\end{center}

\subsection{Perpendicular Lines and the Origin of the Slope Criterion}

Two intersecting lines are called \emph{perpendicular} when the angle between
them is a right angle, that is, an angle of measure
\[
\frac{\pi}{2}.
\]
This is again a geometric definition.  Merely having one common point is not
enough: most pairs of intersecting lines do not form a right angle.

To connect this definition with slope, consider a non-vertical line and let
$\alpha$ be the angle through which the positive $x$-axis must be rotated
counterclockwise to point in the direction of the line.  A right triangle drawn
along the line has horizontal side $\Delta x$ and vertical side $\Delta y$.
By the definition of tangent,
\[
\tan\alpha=\frac{\Delta y}{\Delta x}=m.
\]
Thus slope is not an unrelated coefficient: it is the tangent of the geometric
angle of inclination.

Now let two non-vertical, non-horizontal lines have inclination angles
$\alpha_1$ and $\alpha_2$.  They are perpendicular precisely when their
directions differ by a right angle.  We may therefore write
\[
\alpha_2=\alpha_1+\frac{\pi}{2}
\]
up to reversing the direction of a line, which does not change the line itself.
Using
\[
\tan\left(\alpha+\frac{\pi}{2}\right)
=-\frac{1}{\tan\alpha},
\]
we obtain
\[
m_2
=\tan\alpha_2
=-\frac{1}{\tan\alpha_1}
=-\frac{1}{m_1}.
\]
Hence
\[
\boxed{m_1m_2=-1.}
\]
This formula has now been derived from the geometric definition of
perpendicularity.  It is a test for a right angle, not the meaning of a right
angle.

The excluded case is also geometrically clear.  A horizontal line has
inclination angle $0$ and slope $0$.  Rotating its direction by $\pi/2$ gives a
vertical line, whose slope is undefined.  Thus horizontal and vertical lines
are perpendicular, even though the product-of-slopes formula cannot be written
for them.

\begin{center}
\begin{tikzpicture}[scale=0.85]
    \draw[axis] (-3.8,0) -- (4.0,0) node[right] {$x$};
    \draw[axis] (0,-3.4) -- (0,3.6) node[above] {$y$};
    \draw[subset,domain=-3:3,samples=2] plot (\x,{0.5*\x});
    \draw[subset,domain=-1.6:1.6,samples=2] plot (\x,{-2*\x});
    \draw[black] (0.22,0.11) -- (0.11,0.33) -- (-0.11,0.22);
    \node[subset label,right] at (2.7,1.35) {$L_1$};
    \node[subset label,left] at (-1.25,2.5) {$L_2$};
\end{tikzpicture}
\end{center}

The logical order is therefore:
\[
\text{geometric definition}
\longrightarrow
\text{coordinate analysis}
\longrightarrow
\text{algebraic criterion}.
\]
A criterion is useful only after we know what it recognises and why it works.
